% Public text-only snapshot of the instructor's Module 12 diffusion slide deck.
% Upstream: /Users/setup/Library/CloudStorage/Box-Box/Teaching/6050/LaTeX/Module 12/12.2S Diffusion.tex
% Upstream SHA-256: eeae4146e886e31a534682efc9952fac4ac508932f59a3f0993f7d928a9401e8
% Local preamble dependencies, TikZ source, executable code, external assets,
% third-party-derived material without a clear provenance boundary, product-current
% surveys, hardware claims, and assignment-specific material are omitted. The book
% independently checks and derives every published claim, implementation, experiment,
% and figure.
\documentclass{article}
\usepackage{amsmath}
\usepackage{amssymb}
\usepackage[margin=1in]{geometry}
\usepackage{hyperref}

\title{Lecture 12.2: Generative AI and Diffusion Models}
\author{Heman Shakeri}
\date{}

\begin{document}
\maketitle

\section*{Public Snapshot Boundary}

This standalone, text-only snapshot preserves the instructor-authored lecture sequence
and selected equations. Diagram source, implementation details, external assets,
release-specific examples, and empirical prescriptions have been removed. The
remaining notes document provenance; they are not a substitute for checking the
primary literature.

\section{Self-Supervision as a Foundation}

In supervised learning, an external label supplies the prediction target. In
self-supervised learning, the target is constructed from the observation itself. A
model can predict a hidden token from visible tokens, a future element from a prefix,
or a clean observation from a corrupted one. In each case, the data supplies both the
input and the target.

For a sequence $\mathbf{x}=(x_1,\ldots,x_T)$, an autoregressive model uses the chain
rule factorization

\[
p(\mathbf{x})=\prod_{t=1}^{T}p(x_t\mid x_1,\ldots,x_{t-1}).
\]

This formulation turns generation into a sequence of conditional predictions. The
diffusion formulation uses a different self-supervised target: it corrupts a complete
observation and learns information needed to reverse that corruption.

% [The original visual comparison of generation strategies is omitted.]
% [The original claim that independently injected noise by itself resolves conditional
% averaging is omitted; that conclusion depends on the probabilistic objective,
% parameterization, and sampling procedure.]

\section{The Forward Diffusion Process}

Let $\mathbf{x}_0\sim p_{\mathrm{data}}$. A prescribed Markov chain adds Gaussian
noise in small steps:

\[
q(\mathbf{x}_t\mid\mathbf{x}_{t-1})
=\mathcal{N}\!\left(
\mathbf{x}_t;
\sqrt{1-\beta_t}\,\mathbf{x}_{t-1},
\beta_t I
\right).
\]

With $\alpha_t=1-\beta_t$ and
$\bar\alpha_t=\prod_{s=1}^{t}\alpha_s$, the marginal at any timestep is

\[
q(\mathbf{x}_t\mid\mathbf{x}_0)
=\mathcal{N}\!\left(
\mathbf{x}_t;
\sqrt{\bar\alpha_t}\,\mathbf{x}_0,
(1-\bar\alpha_t)I
\right).
\]

Thus a training pair at time $t$ can be constructed directly:

\[
\mathbf{x}_t
=\sqrt{\bar\alpha_t}\,\mathbf{x}_0
+\sqrt{1-\bar\alpha_t}\,\boldsymbol{\epsilon},
\qquad
\boldsymbol{\epsilon}\sim\mathcal{N}(0,I).
\]

\section{Scores, Noise Prediction, and the Reverse Process}

For a density $p_t$ at noise level $t$, the score is

\[
s_t(\mathbf{x})=\nabla_{\mathbf{x}}\log p_t(\mathbf{x}).
\]

For the Gaussian corruption above, the lecture states the conditional-noise identity

\[
\nabla_{\mathbf{x}_t}\log p_t(\mathbf{x}_t)
=-\frac{1}{\sqrt{1-\bar\alpha_t}}
\mathbb{E}[\boldsymbol{\epsilon}\mid\mathbf{x}_t].
\]

This connects denoising to score estimation. A network
$\boldsymbol{\epsilon}_{\theta}(\mathbf{x}_t,t)$ can be trained with

\[
\mathcal{L}
=\mathbb{E}_{t,\mathbf{x}_0,\boldsymbol{\epsilon}}
\left[
\left\|
\boldsymbol{\epsilon}
-\boldsymbol{\epsilon}_{\theta}(\mathbf{x}_t,t)
\right\|_2^2
\right].
\]

A conceptual training step samples an observation, a timestep, and Gaussian noise;
constructs $\mathbf{x}_t$ with the closed form above; and updates the denoiser toward
the sampled noise. Generation starts with
$\mathbf{x}_T\sim\mathcal{N}(0,I)$ and repeatedly applies a reverse transition. Under
one DDPM parameterization, its mean is

\[
\boldsymbol{\mu}_{\theta}(\mathbf{x}_t,t)
=\frac{1}{\sqrt{\alpha_t}}
\left(
\mathbf{x}_t
-\frac{1-\alpha_t}{\sqrt{1-\bar\alpha_t}}
\boldsymbol{\epsilon}_{\theta}(\mathbf{x}_t,t)
\right).
\]

The variance choice and any stochastic term must be specified alongside this mean;
the equation alone is not a complete sampler.

\section{Provenance Boundary for the Creativity Segment}

The original deck's ``creativity as structured failure'' segment is omitted in its
entirety. Its black/white examples, patch-mosaic argument, locality and equivariance
discussion, coarse-to-fine account, and spatial-inconsistency claims closely track a
primary-paper result without a sufficiently clear independent provenance boundary.
Consult Kamb and Ganguli,
\href{https://proceedings.mlr.press/v267/kamb25a.html}
{``An Analytic Theory of Creativity in Convolutional Diffusion Models''}
(ICML 2025) directly. If the topic enters the book, it must be independently explained,
scoped, and verified from that primary paper.

\section{A Multi-Resolution Denoiser}

The lecture returns to an encoder--decoder with skip connections as a useful denoiser
for image-shaped data. The noisy observation and timestep are inputs; downsampling
builds multi-scale features, while upsampling and skip connections return a prediction
at the original spatial resolution. The architecture predicts a field with the same
shape as the corruption target.

% [The original U-Net diagrams and architecture-specific prescriptions are omitted.]

\section{Latent Diffusion}

Latent diffusion moves the corruption and reverse processes from observation space to
a learned representation. With an encoder $\mathcal{E}$ and decoder $\mathcal{D}$,

\[
\mathbf{z}_0=\mathcal{E}(\mathbf{x}),
\qquad
\hat{\mathbf{x}}=\mathcal{D}(\mathbf{z}_0).
\]

The denoiser acts on noisy versions of $\mathbf{z}_0$, and a generated latent is
decoded afterward. This changes the dimensionality and geometry of the diffusion
state. It does not, without a specified model and benchmark, establish a particular
speed, memory, or quality result.

\section{Conditional Generation}

Cross-attention can connect a noisy representation to a conditioning sequence:

\[
\operatorname{Attention}(Q,K,V)
=\operatorname{softmax}\!\left(\frac{QK^{\top}}{\sqrt d}\right)V.
\]

The noisy representation supplies queries, while the conditioning representation
supplies keys and values. Classifier-free guidance combines unconditional and
conditional noise predictions. Under the convention used in this deck,

\[
\tilde{\boldsymbol{\epsilon}}
=\boldsymbol{\epsilon}_{\theta}(\mathbf{x}_t,t,\varnothing)
+s\left[
\boldsymbol{\epsilon}_{\theta}(\mathbf{x}_t,t,\mathbf{c})
-\boldsymbol{\epsilon}_{\theta}(\mathbf{x}_t,t,\varnothing)
\right].
\]

The scale convention must be stated explicitly because equivalent descriptions use
differently shifted coefficients. Increasing the coefficient changes the trade-off
between conditioning and the unguided model; no universal setting follows from the
formula.

% [The named-product pipeline, release timeline, sampling-speed survey, numerical
% metric table, hardware prescriptions, and all executable examples are omitted.]

\section{Primary Reading for the Retained Equations}

\begin{itemize}
\item Ho, Jain, and Abbeel,
\href{https://arxiv.org/abs/2006.11239}
{``Denoising Diffusion Probabilistic Models''} (2020).
\item Song et al.,
\href{https://openreview.net/forum?id=PxTIG12RRHS}
{``Score-Based Generative Modeling through Stochastic Differential Equations''}
(2021).
\item Rombach et al.,
\href{https://openaccess.thecvf.com/content/CVPR2022/html/Rombach_High-Resolution_Image_Synthesis_With_Latent_Diffusion_Models_CVPR_2022_paper.html}
{``High-Resolution Image Synthesis with Latent Diffusion Models''} (2022).
\item Ho and Salimans,
\href{https://arxiv.org/abs/2207.12598}
{``Classifier-Free Diffusion Guidance''} (2022).
\end{itemize}

\section{Lecture Through-Line}

The lecture's progression is: construct targets from data, define a known corruption
process, learn a time-conditioned denoiser, run the learned process in reverse, then
add representation learning and conditioning. The retained equations record that
conceptual spine while leaving implementations and empirical comparisons to be rebuilt
and tested independently.

\end{document}
